\chapter{极限}

\section{滤子的极限}

记号约定:$X$是拓扑空间,所装备的拓扑记作$\tau$,$a\in X$,对每个$x\in X$,它在$X$中的邻域滤子记作$\mathfrak{B}\left(x\right)$.

\begin{definition}{滤子和滤子基的收敛}{}
	\begin{enumerate}
		\item 设$\mathfrak{F}$是$X$上的滤子.若$\mathfrak{B}\left(a\right)\subseteq\mathfrak{F}$,则称$a$是$\mathfrak{F}$在$X$中的极限点,$\mathfrak{F}$在$X$中(按$\tau$)下收敛到$a$.
		\item 设$\mathfrak{B}$是$X$上的滤子基.若$\mathfrak{B}$生成的滤子收敛到$a$,则称$a$是$\mathfrak{B}$在$X$中的极限点,$\mathfrak{B}$在$X$中(按$\tau$)下收敛到$a$.
	\end{enumerate}
\end{definition}

\begin{proposition}{滤子基收敛的等价条件}{}
	设$\mathfrak{B}$是$X$上的滤子基,则$\mathfrak{B}$收敛到$a$ $\iff$ $a$在$X$中的每个邻域都被$\mathfrak{B}$的某个元素所包含.
\end{proposition}

\begin{remark}{}{}
	\begin{enumerate}
		\item 用邻域基来描述该命题:取定$a$的某个邻域基,则
		\begin{center}
			$\mathfrak{B}$收敛到$a$ $\iff$该邻域基的每个元素都被$\mathfrak{B}$的某个元素所包含.
		\end{center}
		\item 若$X$上的滤子$\mathfrak{F}$收敛到$a$,则$X$上每个比$\mathfrak{F}$细的滤子都收敛到$a$.进一步,把$\tau$换成更粗的拓扑,相应地替换$a$的邻域滤子,则$\mathfrak{F}$在新拓扑下仍收敛.
		
		因此,\underline{一个拓扑越细,那么在此拓扑下收敛的滤子就越少}.特别地,在离散拓扑下收敛的滤子只有邻域滤子,因为它们是平凡的超滤子.
	\end{enumerate}
\end{remark}

\begin{proposition}{收敛到同一点的滤子的交仍然收敛}{}
	设$\Phi$是$X$上一些收敛到$a$的滤子,则$\bigcap\limits_{\mathfrak{F}\in\Phi}\mathfrak{F}$也收敛到$a$.
\end{proposition}

\begin{proposition}{收敛性的超滤子刻画}{}
	设$\mathfrak{F}$是$X$上的滤子,则$\mathfrak{F}$收敛到$a$ $\iff$ $X$上每个比$\mathfrak{F}$细的超滤子都收敛到$a$.
\end{proposition}

\begin{remark}{}{}
	一般而言,在拓扑空间中,一个滤子(或滤子基)有可能收敛到不同的点.
\end{remark}











































































